% ============================================================
\chapter{Conclusion, Future Work, and Industry Approaches}
% ============================================================
\label{ch:conclusion}


\vspace{0.5em}
\noindent\rule{\textwidth}{0.4pt}
\noindent\textbf{Declaration of Contributions:}
This chapter was compiled and authored by \textit{Vishesh Gupta}, who
interpreted the fairness metrics, analysed experimental results, and connected
the technical findings with real-world AI systems deployed by industry leaders.
\textit{Nitin Agiwal} contributed the concluding synthesis of bias and fairness
across the machine-learning pipeline. All team members cross-checked factual claims and references for accuracy.
\noindent\rule{\textwidth}{0.4pt}
\vspace{1em}

\bigskip

This chapter consolidates the key findings of the report, outlines directions
for future investigation, and situates the research within the broader landscape
of how major AI organisations currently address algorithmic fairness.
We begin with a synthesis of our conclusions, followed by concrete open
problems, and close with a comparative survey of fairness approaches adopted by
OpenAI, Google, and IBM.

% ============================================================
\section{Conclusion}
% ============================================================

Bias in AI is a systemic problem deeply embedded across the entire
machine-learning pipeline, from data collection and model training through to
evaluation.
It arises due to historical inequalities reflected in data, underrepresentation
of certain demographic groups, and human subjectivity in labelling and feature
selection~\cite{propublica_compas}.
Real-world case studies such as the COMPAS recidivism algorithm and Amazon's
hiring system demonstrate that such biases can directly harm marginalised
communities and perpetuate existing social inequalities at scale.

Traditional evaluation metrics such as accuracy are insufficient for assessing
fairness: models can achieve high overall performance while producing
disproportionately harmful outcomes for specific demographic groups.
This work examined the FD-VAE architecture~\cite{park2020readme} as a
principled approach to fair representation learning.
FD-VAE addresses the fundamental flaw in prior methods, most notably
FFVAE, by introducing a three-way disentanglement of the latent space into
\emph{target} ($z_t$), \emph{protected} ($z_p$), and \emph{mutual} ($z_m$)
attribute components, ensuring that the shared predictive information is
exploited without leaking demographic signals into downstream classifiers.

The incorporation of Conditional Value-at-Risk (CVaR) loss~\cite{ju2024fairness}
further strengthens robustness by focusing training on worst-performing
subgroups rather than on average-case behaviour, encouraging fairness even in
the absence of explicit group labels.
The proposed pipeline extends these ideas to deepfake detection by combining
Fractional Fourier Transform (FrFT) preprocessing~\cite{ozaktas2001frft},
latent consistency enforcement, adversarial debiasing, and CVaR fairness loss.
Together, these components aim to build a detection system that is not only
accurate but equitable across demographic groups.

The ultimate goal is to develop AI systems that are \textbf{trustworthy},
\textbf{socially responsible}, and \textbf{fair by design} rather than as an
afterthought.

% ============================================================
\section{Future Work}
% ============================================================

Several directions remain open for further investigation.

\paragraph{Empirical Validation.}
The proposed FrFT-based pipeline has been designed at an architectural level;
empirical validation on standard benchmarks such as CelebA and UTK Face
datasets is needed to quantitatively confirm its fairness--accuracy trade-offs
against existing baselines including FFVAE and FD-VAE~\cite{park2020readme}.

\paragraph{Multi-Attribute Fairness.}
The current approach assumes binary protected attributes such as gender.
Extending the framework to handle multiple intersecting sensitive attributes
simultaneously, race, age, and gender combined, presents a non-trivial
challenge that warrants dedicated study.
Intersectional fairness metrics that jointly penalise subgroup disparities
would be a natural first step.

\paragraph{CVaR Hyperparameter Scheduling.}
The CVaR parameter $\alpha$ critically governs the fairness--accuracy
trade-off~\cite{ju2024fairness}.
A systematic sensitivity analysis and an adaptive scheduling strategy for
$\alpha$ during training could improve generalisation across datasets with
varying demographic distributions.

\paragraph{Training Stability of Adversarial Objectives.}
While adversarial debiasing is effective, it introduces the training
instability common to minimax objectives.
Exploring alternative decorrelation strategies, such as kernel-based
independence tests or information-theoretic regularisation, may yield more
stable and scalable training dynamics.

\paragraph{Broader Application Domains.}
The applicability of this pipeline to high-stakes domains beyond deepfake
detection, including medical imaging, credit scoring, and hiring
systems, represents a promising and socially impactful direction.
Each domain introduces unique distributional challenges (e.g.\ class imbalance
in medical data, legally protected attributes in credit scoring) that would
require careful adaptation of the proposed fairness constraints.

% ============================================================
\section{How AI Industry Leaders Address Fairness}
% ============================================================

Beyond academic research, major technology companies have developed their own
fairness frameworks and tools.
These approaches span training-time interventions, post-processing audits, and
open-source toolkits.
Three leading efforts are examined below.

% , , , , , , , , , , , , , , , , , , , , 
\subsection{OpenAI ,  Reinforcement Learning with Human Feedback (RLHF)}
% , , , , , , , , , , , , , , , , , , , , 

OpenAI employs Reinforcement Learning with Human Feedback
(RLHF)~\cite{ouyang2022rlhf} to align model behaviour with human values,
including fairness.
The key techniques include:

\begin{itemize}[leftmargin=2em]
    \item \textbf{Human reviewers} who guide the model toward safe, unbiased
          responses through preference-based feedback.
    \item \textbf{Bias evaluation benchmarks} that test model outputs across
          demographic groups before deployment.
    \item \textbf{Content moderation layers} that filter outputs in real time
          to prevent discriminatory or toxic content.
\end{itemize}

The primary goal of RLHF in this context is to ensure that language models
avoid generating discriminatory responses or perpetuating harmful stereotypes.
Human oversight is central to this approach, making it effective but also
resource-intensive and dependent on the diversity of human reviewers.

% , , , , , , , , , , , , , , , , , , , , 
\subsection{Google ,  Responsible AI Policies}
% , , , , , , , , , , , , , , , , , , , , 

Google has embedded fairness considerations into its model development
lifecycle through its Responsible AI framework, applied prominently in the
Gemini family of models~\cite{gemini2023}.
The key techniques are:

\begin{itemize}[leftmargin=2em]
    \item \textbf{Safety tuning during training}, where dataset composition and
          sampling strategies are adjusted to reduce demographic disparities
          before training begins.
    \item \textbf{Dataset filtering} to remove or rebalance data that may
          encode historical inequalities.
    \item \textbf{Fairness checks} that explicitly test for gender-, race-, and
          culture-based disparities in model outputs.
\end{itemize}

Google's approach addresses bias at the source, during data preparation and
model training, and is generally more effective than post-hoc correction.
However, ensuring comprehensive coverage across all demographic dimensions
remains an ongoing challenge.

% , , , , , , , , , , , , , , , , , , , , 
\subsection{IBM ,  AI Fairness 360 (AIF360)}
% , , , , , , , , , , , , , , , , , , , , 

IBM has released AI Fairness 360 (AIF360)~\cite{aif360}, an open-source Python
toolkit designed to help practitioners detect and mitigate bias across the full
machine-learning pipeline.
The toolkit provides:

\begin{itemize}[leftmargin=2em]
    \item \textbf{Bias detection algorithms} that quantify disparate impact,
          statistical parity, and other fairness metrics.
    \item \textbf{Bias mitigation algorithms} categorised by intervention stage:
    \begin{itemize}[leftmargin=1.5em]
        \item \textit{Pre-processing:} Reweighting and disparate impact remover
              applied to training data to reduce bias at its source, ensuring
              that the model learns from a more balanced representation of the
              data.
        \item \textit{In-processing:} Adversarial debiasing and prejudice
              remover applied during model training, actively guiding the
              learning algorithm to reduce bias while it is being trained.
        \item \textit{Post-processing:} Calibrated equalized odds and reject
              option classification applied to model outputs, adjusting final
              predictions to ensure more equitable outcomes across different
              groups.
    \end{itemize}
\end{itemize}

AIF360's open-source nature makes it widely accessible and adaptable across
industries.
It supports a wide range of fairness metrics and can be integrated into
existing machine-learning workflows with minimal overhead.

% , , , , , , , , , , , , , , , , , , , , 
\subsection{Comparative Summary}
% , , , , , , , , , , , , , , , , , , , , 

Table~\ref{tab:industry_comparison} summarises the three approaches along the
dimensions of primary strategy, key techniques, and the types of harm each is
designed to avoid.

\begin{table}[H]
\centering
\renewcommand{\arraystretch}{1.35}
\begin{tabularx}{\textwidth}{lXXX}
\toprule
\textbf{Company} & \textbf{Approach} & \textbf{Techniques} & \textbf{What It Avoids} \\
\midrule
\textbf{OpenAI} \newline (ChatGPT / GPT-4) &
    RLHF (Reinforcement Learning with Human Feedback)~\cite{ouyang2022rlhf} &
    Human reviewers guide safe responses; bias evaluation benchmarks; content
    moderation layers &
    Toxic outputs; discriminatory responses \\
\addlinespace
\textbf{Google} \newline (Gemini) &
    Responsible AI Policies~\cite{gemini2023} &
    Safety tuning during training; dataset filtering &
    Gender-, race-, and culture-based disparities (via fairness checks) \\
\addlinespace
\textbf{IBM} \newline (AIF360) &
    Open-source Fairness Toolkit~\cite{aif360} &
    Bias detection; bias mitigation algorithms at pre-, in-, and
    post-processing stages &
    Algorithmic bias across protected groups \\
\bottomrule
\end{tabularx}
\caption{Comparison of industry approaches to AI fairness.}
\label{tab:industry_comparison}
\end{table}